\paragraph{CSE545 Software Security - Reflector Project}

\subparagraph{Deweaponized Reflector Demo (Educational)}

\begin{framed}
\textbf{Important}\\
This page documents an educational packet-reflecting demo created for a software security course.

\textbf{Key functionality has been intentionally obscured/disabled for safety:}

\begin{itemize}
\item \textbf{Live sniffing} is disabled by default.
\item \textbf{Packet injection} (\texttt{sendp}) is replaced by a safe logger that does not transmit frames.
\end{itemize}

This preserves learning value (protocol layers, field meanings, control flow) while preventing use as a live network spoofing/reflector tool.
\end{framed}


\bigskip
\centerline{\rule{13cm}{0.4pt}}
\bigskip

\subparagraph{Threat-model framing (defensive lens)}

\begin{framed}
\textbf{Why this matters in security engineering}\\
This demo illustrates \textit{why} certain network assumptions fail:

\begin{itemize}
\item \textbf{ARP has no built-in authentication} (common in LAN environments).
\item Stateless protocols (e.g., \textbf{UDP/ICMP}) are easier to spoof/reflect than \textbf{TCP} (which is stateful).
\item Forged packets must account for \textbf{checksums}, \textbf{headers}, and \textbf{identity mismatches} (L2 vs L3).
\end{itemize}
\end{framed}

\begin{framed}
\textbf{Warning}\\
Do not run ``live-capture + live-injection'' traffic-manipulation tooling on networks you do not explicitly own and control.
Even in a lab, keep experiments isolated (e.g., host-only / private virtual networks).
This publication keeps the code \textit{non-operational} by design.
\end{framed}


\bigskip
\centerline{\rule{13cm}{0.4pt}}
\bigskip

\subparagraph{What this code demonstrates (and what it does not)}

\begin{framed}
\textbf{Demonstrates (safe):}\\
\begin{itemize}
\item Layer identification (\texttt{Ether/IP/ARP/ICMP/TCP/UDP})
\item Packet-field extraction (src/dst MACs, src/dst IPs, L4 fields)
\item ``Would-have-built'' forged packet objects \textit{without sending}
\end{itemize}

\textbf{Does not demonstrate:}

\begin{itemize}
\item How to deploy or run a real spoofing/reflector tool
\item How to optimize or ``make it work'' on a real network
\end{itemize}
\end{framed}


\bigskip
\centerline{\rule{13cm}{0.4pt}}
\bigskip

\subparagraph{Walkthrough: protocol anatomy}

\subparagraph{Layer 2: Ethernet (L2)}

\subparagraph{Layer 3: IP (L3)}

\subparagraph{ARP}

\subparagraph{ICMP}

\subparagraph{UDP vs TCP}


\bigskip
\centerline{\rule{13cm}{0.4pt}}
\bigskip

\subparagraph{Deweaponized code listing}

\begin{framed}
\textbf{Publication safety changes}\\
The code below is deliberately \textbf{non-operational} as a live tool:

\begin{itemize}
\item \texttt{scapy.sendp(...)} is replaced by \texttt{safe\_sendp(...)} which \textbf{never transmits}.
\item Live sniffing is disabled unless processing an \textbf{offline PCAP} via \texttt{--pcap}.
\end{itemize}
\end{framed}


\bigskip
\centerline{\rule{13cm}{0.4pt}}
\bigskip

\subparagraph{Detailed explanation of the deweaponized script}

\begin{framed}
\textbf{Important}\\
This explanation describes the \textbf{published, deweaponized} version.

Two capabilities that make traffic-manipulation tools dangerous are \textbf{intentionally disabled}:

\begin{enumerate}
\item \textbf{Live packet capture} (disabled unless offline \texttt{--pcap} is provided)
\item \textbf{Packet injection} (all sends are routed through \texttt{safe\_sendp()} which does not transmit)
\end{enumerate}

The remaining logic is preserved to teach:

\begin{itemize}
\item protocol layering (Ether/IP/ARP/ICMP/TCP/UDP),
\item field extraction,
\item and \textit{how a forged packet would be constructed in theory}.
\end{itemize}
\end{framed}


\bigskip
\centerline{\rule{13cm}{0.4pt}}
\bigskip

\subparagraph{1) ``Safety switches'': making the script non-operational on a real network}

The script introduces two global flags:

\begin{verbatim}
DISABLE_LIVE_SNIFFING = True
DISABLE_PACKET_INJECTION = True
\end{verbatim}

These are the most important lines in the publication-safe version.

\begin{itemize}
\item \textbf{\texttt{DISABLE\_LIVE\_SNIFFING}}

\begin{itemize}
\item When \texttt{True}, the script will \textbf{not} call \texttt{scapy.sniff(...)}.
\item Instead, it prints a message and exits unless you provide \texttt{--pcap}.
\end{itemize}


\item \textbf{\texttt{DISABLE\_PACKET\_INJECTION}}

\begin{itemize}
\item When \texttt{True}, the script will \textbf{not} call \texttt{scapy.sendp(...)} even if a packet is ``constructed''.
\item Instead, all ``send'' actions are routed to \texttt{safe\_sendp(...)}, which \textbf{only logs} a summary.
\end{itemize}
\end{itemize}

This means the code can be studied and executed for learning, but it does \textbf{not} function as a live network manipulator.

\begin{framed}
\textbf{Why disabling both capture and injection matters}\\
A live traffic tool needs \textbf{both} of the following to have real-world impact:

\begin{itemize}
\item a way to \textbf{observe} traffic (capture),
\item a way to \textbf{affect} traffic (injection).
\end{itemize}

Disabling injection alone prevents harm, but disabling capture as well further reduces the chance that the code can be repurposed by copy/paste.
\end{framed}


\bigskip
\centerline{\rule{13cm}{0.4pt}}
\bigskip

\subparagraph{2) \texttt{safe\_sendp(...)}: a safe stand-in for Scapy injection}

In Scapy:

\begin{itemize}
\item \texttt{send()} transmits at L3 (IP)
\item \texttt{sendp()} transmits at L2 (Ethernet frames)
\end{itemize}

The original script used \textbf{\texttt{sendp()}} because it wanted to control Ethernet headers directly.

The deweaponized script replaces this with:

\begin{verbatim}
def safe_sendp(frame, iface=None, verbose=False):
    print("[SAFE MODE] Blocked packet injection. Would have sent:", frame.summary())
\end{verbatim}

Key points:

\begin{itemize}
\item It \textbf{never transmits} a frame.
\item It avoids dumping raw bytes.
\item It uses \texttt{frame.summary()} which produces a compact ``layer chain'' string like:

\begin{itemize}
\item \texttt{Ether / IP / ICMP 192.168... \textgreater  192.168...}
\item \texttt{Ether / ARP is-at ...}
\end{itemize}
\end{itemize}

This preserves the learning value (``what would have been sent?'') without creating real network effects.


\bigskip
\centerline{\rule{13cm}{0.4pt}}
\bigskip

\subparagraph{3) Command-line arguments: controlling identities for an educational model}

The script defines five core parameters:

\begin{itemize}
\item \texttt{--interface}
\item \texttt{--victim-ip}, \texttt{--victim-ethernet}
\item \texttt{--reflector-ip}, \texttt{--reflector-ethernet}
\end{itemize}

These represent \textit{identities} the script uses to explain packet construction.

\begin{framed}
\textbf{Note}\\
Even though the script accepts \texttt{--interface}, the published version does not perform live sniffing by default.
The interface is used for validation and for labeling/logging.
\end{framed}

It also adds a \textbf{safe} option:

\begin{itemize}
\item \texttt{--pcap \textless file.pcap\textgreater }
\end{itemize}

This enables \textbf{offline packet processing}:

\begin{itemize}
\item No network capture
\item No injection
\item Just parsing and demonstrating the control flow
\end{itemize}


\bigskip
\centerline{\rule{13cm}{0.4pt}}
\bigskip

\subparagraph{4) Input validation: why it exists and what it prevents}

The script validates IPs and MACs before proceeding.

\subparagraph{IP validation (\texttt{ipaddress.ip\_address})}

\begin{verbatim}
ipaddress.ip_address(args.victim_ip)
ipaddress.ip_address(args.reflector_ip)
\end{verbatim}

\begin{itemize}
\item Ensures \texttt{victim\_ip} and \texttt{reflector\_ip} are syntactically valid IP literals.
\item Prevents confusing runtime errors later during packet building.
\end{itemize}

\subparagraph{MAC validation (regex)}

\begin{verbatim}
'^([0 -9A -Fa -f]{2}[: -]){5}([0 -9A -Fa -f]{2})$'
\end{verbatim}

\begin{itemize}
\item Enforces 6 octets in hex format (e.g., \texttt{08:00:27:aa:bb:cc}).
\item Rejects malformed strings early.
\end{itemize}

\subparagraph{Interface validation (\texttt{netifaces})}

\begin{verbatim}
if args.interface not in netifaces.interfaces(): ...
\end{verbatim}

\begin{itemize}
\item A basic sanity check.
\item In the published version it's wrapped in a try/except, so readers without \texttt{netifaces} can still run offline PCAP demonstrations.
\end{itemize}

\begin{framed}
\textbf{Defensive engineering point}\\
Even ``educational'' tools benefit from strict input validation:

\begin{itemize}
\item fewer ambiguous errors,
\item safer defaults,
\item clearer demonstration of expected formats.
\end{itemize}
\end{framed}


\bigskip
\centerline{\rule{13cm}{0.4pt}}
\bigskip

\subparagraph{5) Core Scapy concepts used in the script}

The code repeatedly uses three Scapy patterns:

\subparagraph{(A) Layer presence tests}

Examples:

\begin{itemize}
\item \texttt{scapy.ARP in packet}
\item \texttt{packet.haslayer(scapy.IP)}
\item \texttt{packet.haslayer(scapy.TCP)}
\end{itemize}

These determine what protocol is present.

\subparagraph{(B) Field extraction by layer indexing}

Examples:

\begin{itemize}
\item \texttt{packet[scapy.Ether].src} and \texttt{.dst}
\item \texttt{packet[scapy.IP].src} and \texttt{.dst}
\item \texttt{packet[scapy.UDP].sport} and \texttt{.dport}
\end{itemize}

This is Scapy's ``decode and index'' interface.

\subparagraph{(C) Packet construction by layering with \texttt{/}}

Example:

\begin{verbatim}
scapy.Ether(...) / scapy.IP(...) / scapy.UDP(...) / payload
\end{verbatim}

The \texttt{/} operator stacks layers from L2 \rightarrow L3 \rightarrow L4 \rightarrow payload.


\bigskip
\centerline{\rule{13cm}{0.4pt}}
\bigskip

\subparagraph{6) ARP section: demonstrate how a forged reply is structured (without sending)}

Function:

\begin{verbatim}
def spoof_ARP_from_Reflector_to_Attacker(packet):
    ...
\end{verbatim}

\subparagraph{Step-by-step behavior}

\begin{enumerate}
\item Confirm the packet contains ARP:

\begin{verbatim}
if scapy.ARP not in packet: return
\end{verbatim}


\item Only act on ARP \textbf{requests} (\texttt{op == 1}):

\begin{verbatim}
if packet[scapy.ARP].op != 1: return
\end{verbatim}


\item Extract key ARP fields:

\begin{itemize}
\item \texttt{pdst}: ``Which IP is being asked about?''
\item \texttt{psrc}: ``Who is asking?'' (requester IP)
\item \texttt{hwsrc}: requester MAC
\end{itemize}


\item Decide whether the request is ``interesting'':

\begin{itemize}
\item if \texttt{pdst} matches \texttt{victim\_ip} or \texttt{reflector\_ip}
\end{itemize}


\item Construct an ARP \textbf{reply} packet object (for demonstration only):

\begin{itemize}
\item \texttt{op=2} (reply)
\item \texttt{psrc=\textless spoofed\_ip\textgreater } (the IP we claim to own)
\item \texttt{hwsrc=reflector\_ethernet} (the MAC we claim is associated)
\item \texttt{pdst=\textless requester\_ip\textgreater }, \texttt{hwdst=\textless requester\_mac\textgreater }
\end{itemize}


\item Wrap it in an Ethernet frame:

\begin{itemize}
\item \texttt{Ether.src = reflector\_ethernet}
\item \texttt{Ether.dst = requester MAC}
\end{itemize}


\item Log the summary and block transmission via \texttt{safe\_sendp}.
\end{enumerate}

\begin{framed}
\textbf{Educational takeaway}\\
ARP is fundamentally a \textbf{cache-priming protocol} on many LANs.
This demo shows the \textit{shape} of a reply, and why defenders watch for unusual IP\rightarrowMAC claims.
\end{framed}


\bigskip
\centerline{\rule{13cm}{0.4pt}}
\bigskip

\subparagraph{7) ICMP section: demonstrate a ``reflected'' packet shape}

There are two ICMP functions:

\subparagraph{(A) ``Reflector identity'' ICMP}

\begin{verbatim}
send_ICMP_packet_from_Reflector_to_Attacker(packet)
\end{verbatim}

\begin{itemize}
\item Requires Ether + IP + ICMP.
\item Builds a new packet addressed back to the sender:

\begin{itemize}
\item \texttt{Ether.dst} becomes original sender MAC
\item \texttt{IP.dst} becomes original sender IP
\end{itemize}


\item Spoofs the \textit{source IP} as \texttt{reflector\_ip}.
\end{itemize}

It copies:

\begin{itemize}
\item \texttt{type}, \texttt{code}, \texttt{id}, \texttt{seq},
\item and the payload.
\end{itemize}

\subparagraph{(B) ``Victim identity'' echo reply}

\begin{verbatim}
send_ICMP_packet_from_Victim_to_Attacker(packet)
\end{verbatim}

\begin{itemize}
\item Constructs an ICMP \textbf{echo reply} (\texttt{type=0}),
\item but sets \texttt{IP.src = victim\_ip} to demonstrate how ``identity spoofing'' can look at L3.
\end{itemize}

In both cases:

\begin{itemize}
\item The packet object is constructed,
\item The script logs \texttt{summary()},
\item Transmission is blocked.
\end{itemize}


\bigskip
\centerline{\rule{13cm}{0.4pt}}
\bigskip

\subparagraph{8) UDP section: demonstrate stateless reflection}

Function:

\begin{verbatim}
send_UDP_packet(packet)
\end{verbatim}

It:

\begin{enumerate}
\item Verifies Ether + IP + UDP exist.
\item Chooses a spoofed \texttt{source\_ip} based on \texttt{dst\_ip}:

\begin{itemize}
\item This is a didactic ``swap identity'' rule.
\end{itemize}


\item Constructs:

\begin{itemize}
\item Ether(src=reflector MAC, dst=sender MAC)
\item IP(src=spoofed, dst=sender IP)
\item UDP(sport=original sport, dport=original dport)
\end{itemize}


\item Copies Raw payload if present.
\item Logs summary; does not send.
\end{enumerate}

\begin{framed}
\textbf{Note}\\
UDP is used here because it is easy to illustrate:
it does not require maintaining per-connection state the way TCP does.
\end{framed}


\bigskip
\centerline{\rule{13cm}{0.4pt}}
\bigskip

\subparagraph{9) TCP section: demonstrate field copying and checksum logic (without being a TCP engine)}

Function:

\begin{verbatim}
send_TCP_packet(packet)
\end{verbatim}

This is intentionally a \textbf{packet anatomy demo}, not a full TCP endpoint.

It copies:

\begin{itemize}
\item \texttt{sport}, \texttt{dport},
\item \texttt{seq}, \texttt{ack},
\item \texttt{flags},
\item \texttt{options},
\item payload if present.
\end{itemize}

Then it deletes checksums:

\begin{verbatim}
del reflected[TCP].chksum
del reflected[IP].chksum
\end{verbatim}

Why show this?

\begin{itemize}
\item Because checksums depend on header values.
\item If you change IP source/destination or TCP fields, the checksum must match.
\item Scapy recomputes checksums when sending---here we show the idiom as part of protocol correctness, even though the packet is never transmitted.
\end{itemize}

\begin{framed}
\textbf{Educational takeaway (TCP)}\\
TCP is \textbf{stateful}. A real TCP-speaking endpoint must maintain session state,
sequence number expectations, retransmission logic, etc.

This demo shows:

\begin{itemize}
\item where those fields live,
\item how they can be read/copied,
\item and why checksums matter,
\end{itemize}

without implementing an actual TCP state machine.
\end{framed}


\bigskip
\centerline{\rule{13cm}{0.4pt}}
\bigskip

\subparagraph{10) The dispatcher: mapping ``what you saw'' to ``what you would build''}

Function:

\begin{verbatim}
def handle_packet(packet):
\end{verbatim}

This function is the ``traffic classifier''.

\begin{itemize}
\item If it sees ARP \rightarrow call the ARP handler.
\item Else if it sees IP:

\begin{itemize}
\item ICMP \rightarrow call ICMP constructor based on destination
\item TCP \rightarrow call TCP constructor for victim/reflector destinations
\item UDP \rightarrow call UDP constructor for victim/reflector destinations
\end{itemize}
\end{itemize}

This preserves the original program's pedagogical structure:

\begin{itemize}
\item ``Detect protocol''
\item ``Extract fields''
\item ``Construct a response-shaped packet object''
\end{itemize}


\bigskip
\centerline{\rule{13cm}{0.4pt}}
\bigskip

\subparagraph{11) Main: offline PCAP parsing instead of live sniffing}

At the bottom:

\begin{itemize}
\item If \texttt{--pcap} is provided:

\begin{verbatim}
pkts = scapy.rdpcap(args.pcap)
for p in pkts:
    handle_packet(p)
\end{verbatim}
\end{itemize}

This is the safest way to demonstrate the logic: it processes an existing capture file.

\begin{itemize}
\item If \texttt{--pcap} is not provided:

\begin{itemize}
\item the script exits with a safety message because \texttt{DISABLE\_LIVE\_SNIFFING=True}.
\end{itemize}
\end{itemize}

\begin{framed}
\textbf{Important}\\
This design ensures your published code does not become a ``turnkey live tool''.
It remains an offline educational parser and packet-constructor demonstrator.
\end{framed}


\bigskip
\centerline{\rule{13cm}{0.4pt}}
\bigskip

\subparagraph{Practical notes for readers (safe)}

\begin{framed}
\textbf{How to use this safely (offline)}\\
\begin{itemize}
\item Use a PCAP from a controlled lab environment (or a benign sample capture).
\item Run the script with \texttt{--pcap} to observe:

\begin{itemize}
\item protocol classification,
\item extracted fields,
\item constructed packet summaries.
\end{itemize}
\end{itemize}

No traffic is transmitted; the output is purely explanatory.
\end{framed}


\bigskip
\centerline{\rule{13cm}{0.4pt}}
\bigskip

\subparagraph{Summary}

\begin{framed}
\textbf{Summary}\\
This publication-safe script preserves the \textbf{structure} of the original coursework artifact:

\begin{itemize}
\item argument parsing + validation,
\item protocol identification,
\item header-field extraction,
\item packet-layer construction,
\item checksum correctness idioms,
\end{itemize}

while intentionally obscuring the two critical operational capabilities:

\begin{itemize}
\item \textbf{live sniffing} (disabled),
\item \textbf{packet injection} (blocked).
\end{itemize}
\end{framed}


\bigskip
\centerline{\rule{13cm}{0.4pt}}
\bigskip

\subparagraph{Suggested exercises (safe)}

\begin{framed}
\textbf{Safe lab exercises}\\
\begin{enumerate}
\item Capture a small PCAP from a controlled environment (or use a benign sample PCAP).
\item Run the script with \texttt{--pcap} and observe:

\begin{itemize}
\item protocol dispatch decisions (ARP vs ICMP vs TCP vs UDP)
\item which header fields are extracted
\item summaries of the ``constructed (non-sent)'' packets
\end{itemize}


\item In Wireshark, compare:

\begin{itemize}
\item the original packet headers
\item the demo's printed ``constructed packet'' summary
\end{itemize}
\end{enumerate}
\end{framed}


\bigskip
\centerline{\rule{13cm}{0.4pt}}
\bigskip

\subparagraph{Defensive detection notes}